%!TEX root = /../../springer-mono-format/book.tex
%!TEX root = /../../thesis-trier.tex

Here, we introduce the algebraic structures of relevance to us.

\begin{enumerate} 

\IN{semigroup}
\IN{product in semigroup}
\IN{associative operation}
\item A \emph{semigroup} 
is a tuple $(S, \cdot)$ where $S$ is a set and $\cdot \colon S \times S \to S$
a binary operation, called a \emph{product} and
written as $x\cdot y$ (infix notation) %instead of~$\cdot(x,y)$
for $x,y \in S$, that is \emph{associative}, i.e.,
$x \cdot ( y \cdot z ) = (x \cdot y)\cdot z$ for $x,y,z \in S$.
Often, %in case the context makes clear what operation is meant, 
when it is clear from context what operation $S \times S \to S$ is meant, 
we simply call the set~$S$ a semigroup, i.e., omit the tuple notation,
and we write $xy$ instead of~$x \cdot y$, i.e., we
omit the operation symbol.
	
\IN{identity element}
\IN{monoid}  
\item A \emph{monoid} 
 is a tuple $(M, \cdot, 1)$ where $(M, \cdot)$ is a semigroup and $1 \in M$ is 
 an \emph{identity element}, i.e., $1x = x1 = x$ for all $x \in M$.
 The identity element is unique,
 for if $1,1' \in M$ are identity elements, then $1 = 1 \cdot 1' = 1'$.
   
\IN{inverse operation}
\IN{group}
\IN{inverse element}
\item A \emph{group}  
 is a tuple $(G, \cdot, ^{-1}, 1)$ where $(G, \cdot, 1)$ is a monoid
 and $^{-1} \colon G \to G$ is a unary operation, the \emph{inverse operation},
  written as~$g^{-1}$ instead of $^{-1}(g)$ for $g \in G$,
 such that $gg^{-1} = g^{-1}g = 1$
 for all $g \in G$.
 We call $g^{-1}$ the \emph{inverse element} of $g \in G$, which
 is the unique element $h \in G$
 such that $hg = gh = 1$. 

\IN{semiring}
\item A \emph{semiring} 
 is a tuple $(R, +, \cdot, 0, 1)$
 where $(R, +, 0)$ and $(R, \cdot, 1)$ are monoids
 and the following additional equations are true:
 
 \begin{enumerate} 
 \item $x + y = y + x$ (\emph{commutativity}), \IN{commutativity}
 
 \item $x \cdot 0 = 0\cdot x = 0$ (\emph{annihilation law} for $0$ in $R$), \IN{annihilation law}
  
   \IN{left-distributivity}
 \item $x \cdot (y+z) = x\cdot y + x \cdot z$ (\emph{left-distributivity}) and \\ 
   $(y+z) \cdot x = y\cdot x + z\cdot x$ (\emph{right-distributivity}).  \IN{right-distributivity}
 \end{enumerate}
 The operation $+$ is called \emph{addition} in $R$
 and $\cdot$ is called \emph{multiplication} in $R$.
\end{enumerate}

Due to associativity, we omit parentheses in %when writing 
expressions over semigroups.

\begin{example}\label{prelim:ex:algstructures}
\begin{enumerate}
\item $(\Sigma^+, \cdot)$ for some alphabet $\Sigma$ is a semigroup. 

\item $(\Sigma^*, \cdot, \varepsilon)$ for some alphabet $\Sigma$ is a monoid.

\item $(\mathbb N_0, +, 0)$ and $(\mathbb N_0, \cdot, 1)$ are monoids with the usual addition and multiplication 
operations.

\item $([n], +_n, 0)$ and $([n], \cdot_n, 1)$ are monoids for $n \ge 0$
with $x +_n y = \min(x + y, n-1)$ 
and $x \cdot_n y = \min(x\cdot y, n-1)$ \label{def:plusmaln}
for $x,y \in [n]$.

\item The integers $\mathbb Z$ with the usual addition, inversion $x \mapsto -x$
and identity element $0 \in \mathbb Z$ form a group.
	
\item The \emph{Boolean semiring} \IN{Boolean semiring} 
 is the semiring $(\{0,1\}, \land, \lor, 0, 1)$.
 
\item $(\mathbb N_0, +, \cdot, 0, 1)$ and $([n], +_n, \cdot_n, 0, 1)$ for $n \ge 1$
 are semirings. For $n = 2$ the latter semiring equals the Boolean semiring.
\end{enumerate}	
\end{example}

\label{def:subsemigroup}
\IN{morphism!of semigroups}
\IN{isomorphism!of semigroups}
\IN{subsemigroup} 
\IN{product on subsets} 
\IN{commutative semigroup}
\IN{idempotent element}
\index{complex product|see{product on subsets}}
\label{def:product_subsets} 
Let~$(S_1,\cdot_1), (S_2,\cdot_2)$ be two semigroups. 
A \emph{(semigroup) morphism} 
is a map $\varphi \colon S_1 \to S_2$ such that $\applyfunc{xy}{\varphi} = \applyfunc{x}{\varphi} \applyfunc{y}{\varphi}$ 
for all $x,y \in S_1$. We call~$\varphi$ an \emph{isomorphism} if there exists a semigroup morphism
$\psi \colon S_2 \to S_1$ such that $\concat{\varphi}{\psi} = \operatorname{id}_{S_1}$
and $\concat{\psi}{\varphi} = \operatorname{id}_{S_2}$.
If $\varphi$ is bijective, for its inverse function $\varphi^{-1} \colon S_2 \to S_1$
we have $\varphi^{-1}(\varphi(x)\varphi(y))
= \varphi^{-1}(\varphi(xy)) = xy = \varphi^{-1}(\varphi(x)) \varphi^{-1}(\varphi(y))$ for $x,y \in S_1$.
This yields that $\varphi^{-1}$ is a semigroup morphism.
With \cref{prelim:prop:surj-inj-inverses} 
we find that the semigroup isomorphisms are precisely the bijective morphisms.
We say that~$S_1$ is a \emph{subsemigroup}
of~$S_2$, written $S_1 \le_S S_2$,
if $S_1 \subseteq S_2$ and, for all $x,y\in S_1$,
$x \cdot_2 y \in S_1$ and $x \cdot_1 y = x \cdot_2 y$.
For $i > 0$ and $x \in S_1$ we set $x^i = x\cdots x$ ($i$ times).
Let $A,B \subseteq S_1$. The
\emph{product on subsets of~$S_1$} (also called \emph{complex product} in~\cite{Kurzweil04})
is $A \cdot B = \{ ab \mid a \in A, b \in B\}$, written $AB$ for short.
This gives a semigroup~$(\mathcal P(S_1), \cdot)$.
Also, we set $aA = \{a\}A$ and $Aa = A\{a\}$ for $a \in S_1$.
We call~$S_1$ \emph{commutative} if $xy = yx$ for all $x,y \in S_1$.
An element $x \in S_1$ is called \emph{idempotent} if $x^2 = x$.

\begin{lemma} \label{prelim:lem:finite-semigrp-idempotent-power}
	Let $S$ be a finite semigroup. Then for each $s \in S$
	there exists $i > 0$ such that $s^i$
	is idempotent.
\end{lemma}
\begin{proof}
	By finiteness, we have $s^{j + p} = s^j$
	for some $j \ge 0$ and $p > 0$.
	Let $i \ge 0$ be such that $j \le i < j + p$ and $i \equiv 0 \pmod{p}$.
	Then $(s^{i})^2 = s^{i+i} = s^i$.
\end{proof}

\begin{remark} \label{prelim:rem:semigroup-adjoin-identity}
	Let $(S, \cdot)$ be a semigroup.
	By adding some element $1 \notin S$
	and setting $1\cdot 1 = 1$ and $1\cdot s = s\cdot 1 = s$ for all $s \in S$,
	we obtain a monoid $S \cup \{1\}$.	
\end{remark}

\IN{isomorphism!of monoids}
\IN{morphism!of monoids}
\IN{submonoid} 
\IN{group of units}
\label{def:submonoid} 
\label{def:powers-smaller}
\label{def:group of units} 
Let $M_1, M_2$ be two monoids with identity elements $1_1$ and $1_2$, respectively.
We set $x^0 = 1_1$ for $x \in M_1$.
A \emph{monoid morphism} (resp.~\emph{monoid isomorphism}) is a 
semigroup morphism (resp.~semigroup isomorphism) $\varphi \colon M_1 \to M_2$
with $\varphi(1_1) = 1_2$.
When the context is clear, we simply say morphism (resp.~isomorphism) instead
of semigroup or monoid morphism (resp.~isomorphism).
We call~$M_1$ a \emph{submonoid} 
of~$M_2$, written~$M_1 \le_M M_2$, if $M_1$ is a subsemigroup of $M_2$ with $1_2 \in M_1$.
The \emph{group of units} 
of $M_1$ 
is
\[ 
\{ g \in M_1 \mid \exists h \in M_1 : gh = hg = 1 \}.
\] 
The group of units is in fact a group.
Further, for $A \subseteq M_1$,  we set 
$A^0 = \{1\}$, $A^{\le n} = A^0 \cup A \cup A^2 \cup \ldots \cup A^n$ 
for $n \ge 0$, $A^{< 0} = \emptyset$ and $A^{< n} = A^{\le n - 1}$ for $n > 0$. 
For $a \in M_1$ we also write $a^{< n}$ and $a^{\le n}$ for $\{a\}^{< n}$ and $\{a\}^{\le n}$, respectively.

\begin{remark} In monoids, there might exist subsemigroups that are monoids in itself, but not submonoids.	
	For example, in $\{1,s,s^2\}$ with $s^2 = s^3$ and $s1 = 1s = s$, $\{ s^2 \}$ is a subsemigroup
	that is a monoid, but it is not a submonoid.
\end{remark} 

\begin{example}
	Let $M$ be a monoid with identity $1 \in M$.
	Then we 
	obtain a semiring structure~$(\mathcal P(M), \cup, \cdot, \emptyset, 1)$
	from the product on subsets and set union.	
\end{example}

\IN{trivial group}
\IN{subgroup} 
\IN{morphism!of groups}
\label{def:subgroup}
Let $G_1, G_2$ be two groups and $1 \in G_1$ be the identity in $G_1$.
We call~$G_1$ \emph{trivial} 
if $G_1 = \{1\}$.
A \emph{group morphism} $\varphi \colon G_1 \to G_2$
is a monoid morphism such that $\varphi(g^{-1}) = \varphi(g)^{-1}$
for all $g \in G_1$.
Now, for every semigroup morphism $\psi \colon G_1 \to G_2$
we have $\psi(1) = \psi(1\cdot 1) = \psi(1)\psi(1)$,
which implies that $\psi(1)$ is the identity in $G_2$,
and $\psi(g)\psi(g^{-1}) = \psi(gg^{-1}) = \psi(1)$,
so that $\psi(g)^{-1} = \psi(g^{-1})$. Thus, the group morphisms
are precisely the semigroup morphisms
and, technically, 
the definition of group morphisms can be omitted (we included it for conceptual reasons).
We say that~$G_1$ is a \emph{subgroup}
of $G_2$, written $G_1 \le_G G_2$,
if~$G_1$ is a submonoid of $G_2$ and the inverse operation in~$G_1$
agrees with the inverse operation in~$G_2$ for each element of~$G_1$, i.e.,
inversion in~$G_1$ equals the inversion in~$G_2$ restricted to~$G_1$.\footnote{Here, we use the term subgroup only
	w.r.t.~groups. In the literature, for example~\cite{McNaughton71}, 
	subgroups are also defined for semigroups as subsemigroups
	having an identity element (that does not have to be an identity for the whole semigroup)
	and containing inverses w.r.t.~this identity.}
Next, assume $G_1 \le_G G_2$. 

\begin{enumerate}
	
	\index{coset|see{right coset}}
	\IN{left coset} 
	\IN{right coset} 
	\item 
	A subset
	of the form $G_1g$ (resp.~$gG_1$) for some $g \in G_2$
	is called a \emph{right coset} (resp.~\emph{left coset}) 
	of~$G_1$ in $G_2$.
	A right coset is also just called coset.
	The set of right cosets (resp.~left cosets) of $G_1$
	forms a partition of $G_2$. If $G_2$ is finite,
	then each right coset (resp.~left coset) has size $|G_1|$.
	
	\IN{maximal subgroup} 
	\item We say that $G_1$ is a \emph{maximal subgroup} 
	of $G_2$ if $G_1 \ne G_2$ and for every subgroup $H$ of $G_2$ with $G_1 \subseteq H$
	we have $H = G_1$ or $H = G_2$.
	
	\IN{normal subgroup} 
	\item We say that $G_1$ is a \emph{normal subgroup}
	of $G_2$ if $gG_1g^{-1} \subseteq G_1$ for all $g \in G_1$.
	
\end{enumerate}

\label{def:center} 
\index{simple subgroup}
\index{center of group}
\index{abelian group}
\index{order of group}
\index{cyclic group}
A group $G$ is said to be \emph{simple}
if the only normal subgroups of $G$ 
are $\{1\}$ and~$G$.
The \emph{center} of $G$ is 
\[
Z(G) = \{ g \in G \mid \forall h \in G \colon hg = gh \}.
\]
We call~$G$ \emph{cyclic} if $G = \langle \{g\} \rangle_G$ for some $g \in G$.
A commutative group is also often called an \emph{abelian group} in the literature.
For a finite group~$G$,
$|G|$ is also called the \emph{order} of the group.

\index{generated!{submonoid}} 
\index{generated!{subgroup}}
\index{generated!{subsemigroup}}
\index{generator}
\label{def:gensubstructures}

For a subset $T$ of a semigroup (resp.~monoid, group)
we denote by $\langle T \rangle_S$ (resp.~$\langle T \rangle_M$, $\langle T \rangle_G$)
the \emph{subsemigroup} (resp.~\emph{submonoid, subgroup}) \emph{generated by~$T$}, i.e., the 
smallest subsemigroup (resp.~submonoid, subgroup) containing~$T$ in the semigroup (resp.~monoid, group) under consideration.
In this regard, the elements in~$T$ are called the \emph{generators}
of the semigroup~$\langle T \rangle_S$ (resp.~monoid~$\langle T \rangle_M$, group~$\langle T \rangle_G$).
We have $\langle T \rangle_S = \bigcup_{i \ge 1} T^i$,
$\langle T \rangle_M = \bigcup_{i \ge 0} T^i$, $\langle \emptyset \rangle_S = \emptyset$
and $\langle \emptyset \rangle_M = \langle \emptyset \rangle_G = \{1\}$.
Suppose $(G, \cdot, ^{-1}, 1)$ is a finite group.
Then for $g \in G$ there exists $n \ge 1$ such that $g^n = 1$
and so $g^{-1} = g^{n-1}$. 
This implies that the subsemigroups in the semigroup $(G, \cdot)$
are precisely the subgroups of $(G, \cdot, ^{-1}, 1)$
and, for non-empty $T \subseteq G$, we have~$\langle T \rangle_G = \langle T \rangle_M = \langle T \rangle_S$.
So, for finite groups and $T \ne \emptyset$,
we can omit~$S, M$ or $G$  in the index and just write unambiguously~$\langle T \rangle$.

% https://tex.stackexchange.com/questions/690308/im-using-subscripts-how-do-i-to-format-cdot-and-varsubscript-correctly-us
\IN{point-wise product} 
\IN{component-wise product} 
\IN{direct product (semigroup)} 
Let~$S_1, S_2$ be semigroups.
Then $S_1 \times S_2$ with the operation
\[
(s, t)(s', t') = (ss', tt') 
\]
for $(s,t), (s', t') \in S_1 \times S_2$ is called the \emph{direct product} semigroup.
If~$S_1, S_2$ are monoids (resp.~groups) with identities $1 \in S_1$ and $1' \in S_2$, respectively,
then $S_1 \times S_2$ is a monoid with identity $(1,1)$
(resp.~group with identity $(1,1)$ and  inversion $(s,t)^{-1} = (s^{-1}, t^{-1})$).
Let~$I$ be some set. We obtain a 
semigroup $(S_1^I, \odot)$
with the \emph{(point-wise} or \emph{component-wise) product}~$\odot$ 
given by $\applyfunc{x}{(f \odot  g)} = \applyfunc{x}{f} \cdot \applyfunc{x}{g}$
for $f,g \in S_1^I$ and all $x \in I$.
If $S_1$ is a monoid with identity~$1 \in S_1$,
then $S_1^I$ is a monoid where the identity is the constant function $\textbf{1}\in S_1^I$
mapping everything to~$1$.

\IN{component-wise addition}
\begin{example} \label{prelim:ex:component-wise}
	Consider 
	$(\mathbb N_0, +, 0)$, $(\mathbb N_0, \cdot, 1)$,
	$([n], +_n, 0)$ and $([n], \cdot_n, 1)$
	from Example~\ref{prelim:ex:algstructures} and an alphabet $\Sigma$.
	We obtain monoid structures
	$(\mathbb N_0^{\Sigma}, \oplus, \textbf{0})$
	and $(\mathbb N_0^{\Sigma}, \odot, \textbf{1})$
	from $(\mathbb N_0, +, 0)$ and $(\mathbb N_0, \cdot, 1)$, respectively,
	by the construction given above where 
	$\oplus$ is also called the \emph{component-wise addition} on $\mathbb N_0^{\Sigma}$.	
	Similarly, we obtain
	monoid structures 
	$(\mathbb [n]^{\Sigma}, \oplus_n, \textbf{0})$
	and $(\mathbb [n]^{\Sigma}, \odot_n, \textbf{1})$
	on $[n]^{\Sigma}$.
	In the following, unless stated otherwise, 
	we assume these monoid structures on $\mathbb{N}_0^{\Sigma}$ and 
	$[n]^{\Sigma}$, and we write $\oplus$ as $+$ and $\oplus_n$ as $+_n$, 
	since there is no danger of confusion.
\end{example}